% RESEARCH_DRAFT_SHA256: e624b23c93bedbe1f173e5a80084da0fa692fecbae22b4835a61cd3953dea18c
% Layout and source-reference wrappers expanded for the web renderer.
\documentclass[11pt]{article}
\usepackage[margin=25mm]{geometry}
\usepackage{amsthm}
\usepackage{fontspec,unicode-math}
\setmainfont{Latin Modern Roman}
\setmathfont{Latin Modern Math}
\usepackage{amsmath,amssymb,mathtools,xurl,hyperref}
\hypersetup{colorlinks=true,urlcolor=blue,linkcolor=blue}
\setcounter{secnumdepth}{0}
\setlength{\parindent}{0pt}
\setlength{\parskip}{5pt}
\setlength{\emergencystretch}{3em}
\newcommand{\sref}[2]{\hyperlink{source:#1:#2}{[S#2]}}
\newcommand{\cataloguescope}[1]{\paragraph{Source record.}#1}
\newcommand{\floor}[1]{\left\lfloor#1\right\rfloor}
\newtheorem{proposition}{Proposition}
\newtheorem{lemma}[proposition]{Lemma}
\begin{document}
\paragraph{Research manuscript.} OpenAI \texttt{gpt-6-astra}, reasoning \texttt{ultra}; 11 October 2026.
\paragraph{Preserved source body.} The complete catalogue mathematical text and sources follow with layout and source links expanded, including source qualifications and its historical preparation notes. New research begins at the section headed Research attempt.

\section{Exponential Tree Size of Cluster States}\label{statement}

\subsection{Definitions and mathematical statement}

\paragraph{Shared notation.}Standard number systems are $\mathbb N,\mathbb Z,\mathbb Q,\mathbb R,\mathbb C$. Specialized notation and index ranges follow the cited source and the definitions below.


A state tree builds a multiqubit vector from one-qubit vectors using addition and tensor product; tree size counts its leaves, minimizing over such representations. A cluster state is the graph state associated with the specified grid graph.
\paragraph{Record scope and qualifications.}
The catalogue title is retained as metadata. Named certificates, contracts, bounds or frameworks in that title are not asserted as results of the original source extract.
The random-code ensemble, specific states in Shor\textquotesingle s algorithm, and uniformity/probability requirements must be taken from the source.

Source-record limitation: The code-state ensemble and particular states arising in Shor' s algorithm are not identified, so the bundled target cannot be recovered in full.
\paragraph{Statement recorded in the supplied source.}
Do uniform superpositions over random linear codes and the states arising in Shor's algorithm have tree size exponential in $n$; and do cluster states in dimensions two and three likewise have exponential tree size? \href{https://doi.org/10.4171/owr/2005/26}{[S2]}
\paragraph{Source status.}
The supplied catalogue labels this source record open; that metadata was checked against the supplied snapshot on 2026-10-11. The label is preserved for provenance and is not an independent verification of present open status or mathematical correctness. The catalogue literature assessment is dated 2026-08-21.
\subsection{Short English statement}
Prove or refute exponential state-tree size for all the listed code, factoring, and cluster-state families.
\subsection{Sources}
\begin{itemize}
\item[S1] Ulam AI and the UnsolvedMath Contributors, UnsolvedMath: A Curated Collection of Open Mathematics Problems, catalogue entry. Catalogue curation and metadata: CC BY 4.0. The source extract is retained; definitions, English summaries and qualifications are editorial additions.
\url{https://huggingface.co/datasets/ulamai/UnsolvedMath}.
\item[S2] Complexity Theory (2006), Oberwolfach Reports, DOI 10.4171/owr/2005/26, pp. 1431--1504
\url{https://doi.org/10.4171/owr/2005/26}.
\item[S3] Aaronson quantum states paper
\url{https://arxiv.org/abs/quant-ph/0507242}.
\item[S4] Tree-size complexity paper
\url{https://journals.aps.org/pra/abstract/10.1103/PhysRevA.88.012321}.
\item Source-specific Oberwolfach Reports licences, including CC BY-SA 4.0 where applicable, are retained; see THIRD\_PARTY\_NOTICES.md.
\end{itemize}
\paragraph{Preparation and review.}
The original source and UnsolvedMath contributors retain attribution for the source extract. Codex (GPT-6; exact serving revision not exposed) assisted with definitions, English summaries, scope notes and formatting. No new human mathematical authorship or independently verified solution is asserted. Catalogue identity checked against the complete supplied registry and current repository. Distinct upstream records, including repeated questions and broad or damaged extracts, are retained for complete coverage of the source CS category. A repeated record is not a new independent conjecture. Complete mathematical-equivalence and current-literature review has not been performed. Preparation and final review follow contribution prompts C and D, with source gaps disclosed. Structural and rendering checks do not establish a new theorem or open status.
\label{end:statement}

\section{Research attempt}
\subsection{Literature, source context and scope}
The full original catalogue body and source qualifications are preserved above. The relevant primary workshop passages were read in the cached publisher report. Historical claims and subsequent attributed results are distinguished from the mathematical work below. A restricted example, standard lemma or known theorem is not counted as new progress on the entire original target. Missing assumptions and repeated or rhetorical extracts are not silently turned into different solved problems.

\subsection{A separator upper bound and the meaning of exponential}
We test the literal interpretation $2^{\Omega(n)}$ when $n$ denotes the total number of qubits of a rectangular grid cluster state. This is narrower than the bundled code-state, Shor-state and cluster-state question. The 2013 primary tree-size paper discusses a superpolynomial conjecture for the two-dimensional cluster state and local-operator invariance. We use the standard separator expansion below, making no novelty claim. A superpolynomial lower bound is compatible with our upper bound.
For a finite graph $G=(V,E)$, write its graph state as
\[
|G\rangle=2^{-|V|/2}\sum_{x\in\{0,1\}^V}(-1)^{\sum_{uv\in E}x_ux_v}|x\rangle.
\]
Let $S\subseteq V$ have size $s$, and let $G_1,\ldots,G_j$ be the nonempty components of $G-S$. For $z\in\{0,1\}^S$, define $b_v(z)=\sum_{u\in S:uv\in E}z_u$ modulo two for $v\notin S$; let $Z_v$ act by $Z_v|x_v\rangle=(-1)^{x_v}|x_v\rangle$.
\begin{lemma}
There is a state-tree representation satisfying
\[
\operatorname{TS}(|G\rangle)\le 2^s\left(s+\sum_i\operatorname{TS}(|G_i\rangle)\right).
\]
\end{lemma}
\begin{proof}
Partitioning the amplitude sum according to $x|_S=z$ gives the exact identity
\[
|G\rangle=2^{-s/2}\sum_z(-1)^{\sum_{uv\in E(G[S])}z_uz_v}|z\rangle\otimes
\bigotimes_i\left(\prod_{v\in V(G_i)}Z_v^{b_v(z)}\right)|G_i\rangle.
\]
Indeed every edge lies inside $S$, inside one component, or joins $S$ to a component; their phase contributions are exactly the three factors displayed. Normalization agrees because $2^{-s/2}\prod_i2^{-|V(G_i)|/2}=2^{-|V|/2}$. A local $Z_v$ only changes the one-qubit leaves labeled $v$, so it does not increase their count. Tensor products add leaf counts, and the $2^s$ summands may be represented by separate copies of the component trees. Scalar factors can be absorbed into a one-qubit leaf, or distributed through its subtree, without adding leaves. The computational-basis vector on $S$ uses $s$ leaves, proving the bound.
\end{proof}
\begin{proposition}
For a rectangular nearest-neighbor grid in fixed dimension $d\in\{2,3\}$ with $n\ge1$ vertices,
\[
\operatorname{TS}(|G\rangle)\le n\,2^{C_d n^{1-1/d}},\qquad
C_d=\frac1{1-2^{-(1-1/d)}}.
\]
\end{proposition}
\begin{proof}
Set $\alpha=1-1/d$. A middle slice perpendicular to a longest side has size $s=n/\ell\le n^\alpha$, since $\ell\ge n^{1/d}$. Its removal leaves at most two rectangular grids, each of size $n_i\le n/2$, with $\sum_i n_i=n-s$. Use induction on $n$, the case $n=1$ being a one-leaf state. The lemma and the inductive bounds give
\[
\operatorname{TS}(|G\rangle)
\le2^s\left(s+\sum_i n_i2^{C_d n_i^\alpha}\right)
\le n\,2^{n^\alpha+C_d(n/2)^\alpha}
=n\,2^{C_d n^\alpha}.
\]
Empty components contribute nothing, so grids with a side of length one cause no exception.
\end{proof}
\subsection{Verification and remaining gaps}
The separator phase identity was checked on every computational-basis string for a $2$ by $3$ grid and a $2$ by $2$ by $2$ grid. Recursive leaf counts were also computed for small rectangles. Since $\log_2n+C_dn^{1-1/d}=o(n)$, the displayed upper bound excludes $2^{\Omega(n)}$ for this precise qubit-count interpretation. It does not exclude $2^{\Omega(n^c)}$ with a smaller positive $c$, or any superpolynomial lower bound compatible with the estimate. If the source's size parameter denotes a side length or uses exponential in a broader sense, it is not refuted by that phrase alone. The code ensemble and specific Shor states remain unspecified, and no lower bound for the intended cluster-state conjecture is proved. This known-method clarification receives no original-target completion credit.

\section{References and source attribution}
Primary workshop report: \url{https://doi.org/10.4171/owr/2005/26}. Relevant passages read: 2005/26 pp. 1464--1465, the bundled tree-size question and cluster-state dimensions.

\url{https://arxiv.org/html/1303.4843v3}. Read primary tree-size paper, particularly its definition, local-operator invariance, one-dimensional upper-bound discussion and two-dimensional superpolynomial conjecture. The separator calculation here is presented as a standard method, not a new lower-bound result.

All retrievals were checked on 11 October 2026. Original question attribution remains in the preserved Sources section. External theorems are attributed and not claimed as new discoveries or fully verified merely by citation.


\paragraph{Integration review.} Codex (GPT-6; exact serving revision not exposed) reviewed the written arguments and expanded layout and source-reference wrappers. Repository metadata and the production web renderer were checked using Python 3.13 and JavaScript with MathJax 3.2.2. This model-assisted review is not independent human proof verification; no PDF compilation or formal proof checking is asserted.

\section{Completion Estimate}
\noindent\textbf{COMPLETION: 0\%}
No substantiated new progress toward the entire original target is claimed. The worked results and failed route are delimited above; a known result, illustrative model, or auxiliary obstruction is not counted as original-target completion.

\section{Final status and attribution}
\textbf{Attempt status: unresolved.} This describes this research attempt, not a new ruling on the source catalogue's status or on external literature claims.
The mathematics and this self-audit were generated by OpenAI \texttt{gpt-6-astra}, reasoning \texttt{ultra}, on 11 October 2026. Codex (GPT-6; exact serving revision not exposed) orchestrated the batch and contributed formatting/repository review. No human mathematical authorship, independent human proof checking, or formal verification is claimed. The historical preparation paragraphs in the reproduced catalogue body concern that body only.

Contribution prompt A was applied to the complete definition and supplied source record, followed by the adversarial checks of prompt B. Local reproduction files are \texttt{batch016/checks.py}, \texttt{computations.json}, and the per-ID audit JSON. They were not public links at preparation time. Numerical checks are finite and do not prove unrestricted results. No LaTeX engine was found locally; compilation was not performed by this research agent.

\end{document}
